\section*{Discussion}

This study establishes a computational connection between a nanosecond electrical perturbation and an extracellular particle population measured hours later. The principal result is not a single optimized protocol, but a separation of processes that are otherwise confounded in concentration data: electrical exposure and membrane charging, calcium and bioenergetic responses, remodeling and repair, pathway-resolved vesicle release, extracellular disappearance and measurement. This separation allowed the model to reproduce the dominant measured particle-size distribution while making its unresolved dose dependence and weakly constrained particle classes explicit.

The model is consistent with a growing experimental picture in which electrical and calcium signals regulate membrane trafficking at several levels. Direct-current electric fields can increase small-EV secretion through ceramide-dependent budding and increased intraluminal-vesicle formation \cite{huang2026electric}. Although that exposure differs substantially from nsPEF stimulation, it supports the broader premise that membrane polarization can regulate vesiculogenesis. The nsPEF study most closely aligned with our upstream layers found size-dependent permeability, calcium-dependent resealing, osmotic swelling and stiffness loss \cite{kielbik2026selective}. These observations support representing cell geometry and repair capacity as determinants of response rather than treating nominal pulse settings as a transferable dose.

Calcium provides a plausible link between transient permeability and longer-lived vesicle release. Extracellular ATP stimulates both small- and large-EV secretion through P2X7-associated calcium entry in several cancer cell lines \cite{kao2023atp}. Calcium-activated scrambling changes the exposed membrane proteome within minutes \cite{tuomivaara2024slapshot}, and TMEM16E-dependent phosphatidylserine exposure can recruit a compensatory macropinocytic repair response \cite{kim2025tmem16e}. In the present simulations, calcium-dependent remodeling increased budding competence, whereas repair limited progression toward injury-associated output. These relationships are mechanistically plausible but are not yet jointly calibrated in the producer cells used for the particle measurements. Their current value is to define measurements that can distinguish competing explanations.

The fitted observation model reinforces the need to separate biogenesis pathways from operational diameter classes. Distinct extracellular particles can overlap in diameter and composition, and secreted midbody remnants provide an example of a particle class molecularly distinct from both exosomes and plasma-membrane microparticles \cite{peterman2021midbody}. The fitted small- and medium/large-EV medians should consequently be interpreted as effective pathway kernels, not direct identification of exosomes or microvesicles. The near-zero apoptotic-body contribution is likewise not evidence that these particles were absent; the measured window and the available time points supplied little information about that pathway.

The inferred effective half-life summarizes several extracellular processes rather than a molecular degradation constant. Uptake, degradation, adsorption, aggregation, sampling and isolation can all reduce measured concentration. The present data did not allow those processes to be estimated separately. The strong sensitivity to source conversion and extracellular disappearance therefore identifies an experimental ambiguity: an apparent concentration decline can be explained by altered production, altered loss or both. Time-matched sham samples, recovery controls and orthogonal particle assays are needed before assigning the fitted loss to a specific biological process.

The empirical exposure gain improved the fit but marks a limitation of the mechanistic model. Because delivered field, electrode separation, waveform, temperature and complete culture metadata remain unconfirmed, the source labels cannot calibrate membrane voltage or absorbed energy. The gain currently absorbs differences that may arise from dosimetry, viable-cell number, recovery, dilution, instrument response or missing biology. Retaining it as an explicit observation-layer term prevents these effects from being silently incorporated into calcium or vesicle-biogenesis constants. A decisive next test is to replace the gain with measured exposure and early cellular responses and evaluate performance on held-out pulse conditions.

Several additional limitations define the present domain of interpretation. The intracellular equations are reduced and contain provisional constants; acute nanosecond dynamics are compressed into pulse descriptors rather than directly integrated. The model represents a typical producer cell and does not yet propagate population heterogeneity. Repeated-measurement independence, instrument detection limits and the origin of the sham measurements require confirmation. The calibration is in-sample, and the sensitivity designs are screening analyses rather than converged uncertainty quantification. Cargo and potency modules are implemented as interfaces but were not constrained by the present dataset. These limitations preclude claims of biological parameter identification or prospective optimization.

The analysis nevertheless yields a falsifiable experimental program. Measurements of delivered field and temperature define the physical input; early permeability, calcium, phosphatidylserine and viability measurements constrain the transition from exposure to remodeling; pathway-associated markers and electron or imaging-flow methods constrain particle identity; and denser extracellular time courses separate production from disappearance. Because each measurement maps to a defined model interface, negative or discordant results can revise a specific mechanism without discarding the complete pipeline. In this way, the model provides a structured route from bioelectric exposure to EV manufacturing questions while preserving the distinction between simulated mechanism, fitted observation and biological evidence.
